% annotations.tex
% This file should be included after the title pages:
% \thispagestyle{empty}

% ---------------------------------------------------------
% TITLE PAGE
% ---------------------------------------------------------
\begin{center}

{MINISTRY OF SCIENCE AND HIGHER EDUCATION OF THE REPUBLIC OF KAZAKHSTAN}\\
INTERNATIONAL INFORMATION TECHNOLOGY UNIVERSITY\\
FACULTY OF COMPUTER TECHNOLOGIES AND CYBERSECURITY

\vspace{35mm}

Nurshanov D.S.\\
Askarov I.Y.\\
Kassymkhanov T.G.\\
Kaidrakhmetova D.Y.

\vspace{15mm}
\textbf{AI-Guided Prediction of Thermoelastic Wave Propagation in Geological Media}

\vspace{10mm}
{DIPLOMA PROJECT}

\vspace{13mm}
\textbf{Educational Program: 6B06112  Data Science}

\vfill

{Almaty 2026}

\end{center}

\newpage
\thispagestyle{empty}

% ---------------------------------------------------------
% APPROVAL PAGE
% ---------------------------------------------------------
\begin{center}
{MINISTRY OF SCIENCE AND HIGHER EDUCATION OF THE REPUBLIC OF KAZAKHSTAN}\\
INTERNATIONAL INFORMATION TECHNOLOGY UNIVERSITY\\
FACULTY OF COMPUTER TECHNOLOGIES AND CYBERSECURITY\\
DEPARTMENT OF MATHEMATICAL AND COMPUTER MODELING
\end{center}

\vspace{4mm}

\begin{flushright}
\begin{minipage}{6.2cm}
Approved\\
Head of Department,\\
PhD, Associate Professor\\
Abdikalikova Z. T.\\[1mm]
\mbox{\guillemotleft{}\rule{0.8cm}{0.1pt}\guillemotright{}\ \rule{2.2cm}{0.1pt}\ 20\rule{0.4cm}{0.1pt}}
\end{minipage}
\end{flushright}

\vspace{8mm}

\begin{center}
\textbf{DIPLOMA PROJECT}

\vspace{4mm}
\textbf{AI-Guided Prediction of Thermoelastic Wave Propagation in Geological Media}

\vspace{5mm}
\textbf{Educational Program: 6B06112 Data Science}
\end{center}

\vspace{8mm}

\noindent
\begingroup
\setlength{\tabcolsep}{0pt}
\renewcommand{\arraystretch}{1.0}
\begin{tabular*}{\textwidth}{@{\extracolsep{\fill}}l c l@{}}
Done by students: & & \\
Nurshanov D.S. & \rule{3.7cm}{0.1pt} & \mbox{\guillemotleft{}\rule{0.8cm}{0.1pt}\guillemotright{}\ \rule{2.2cm}{0.1pt}\ 20\rule{0.4cm}{0.1pt}} \\
& \small{(signature)} & \\[2mm]
Askarov I.Y. & \rule{3.7cm}{0.1pt} & \mbox{\guillemotleft{}\rule{0.8cm}{0.1pt}\guillemotright{}\ \rule{2.2cm}{0.1pt}\ 20\rule{0.4cm}{0.1pt}} \\
& \small{(signature)} & \\[2mm]
Kassymkhanov T.G. & \rule{3.7cm}{0.1pt} & \mbox{\guillemotleft{}\rule{0.8cm}{0.1pt}\guillemotright{}\ \rule{2.2cm}{0.1pt}\ 20\rule{0.4cm}{0.1pt}} \\
& \small{(signature)} & \\[2mm]
Kaidrakhmetova D.Y. & \rule{3.7cm}{0.1pt} & \mbox{\guillemotleft{}\rule{0.8cm}{0.1pt}\guillemotright{}\ \rule{2.2cm}{0.1pt}\ 20\rule{0.4cm}{0.1pt}} \\
& \small{(signature)} & \\[4mm]
\begin{tabular}[t]{@{}l@{}}
Research advisor:\\
Associate Professor,\\
Candidate of Physical and\\
Mathematical Sciences,\\
PhD in Applied Mathematics\\
Alipova B. N.
\end{tabular}
&
\begin{tabular}[t]{@{}c@{}}
\rule{3.7cm}{0.1pt}\\[-1.5mm]
{\small(signature)}
\end{tabular}
&
\mbox{\guillemotleft{}\rule{0.8cm}{0.1pt}\guillemotright{}\ \rule{2.2cm}{0.1pt}\ 20\rule{0.4cm}{0.1pt}}
\end{tabular*}
\endgroup

\newpage
\thispagestyle{empty}
\enlargethispage{10mm}

% ---------------------------------------------------------
% DIPLOMA PROJECT ASSIGNMENT
% ---------------------------------------------------------
\begin{center}
International Information Technology University\\
Faculty of Computer Technologies and Cybersecurity\\
Department of Mathematical and Computer Modeling
\end{center}

\vspace{9mm}

\noindent Educational Program: 6B06112 -- Data Science

\vspace{4mm}
\noindent Diploma Project Assignment

\vspace{4mm}
\noindent Students\\
\textbf{Nurshanov D.S., Askarov I.Y., Kassymkhanov T.G., Kaidrakhmetova D.Y.}

\vspace{4mm}
\noindent Diploma project topic\\
\textbf{AI-Guided Prediction of Thermoelastic Wave Propagation in Geological Media}

\vspace{4mm}
\noindent Approved by IITU order No. \rule{1.4cm}{0.1pt}\ dated
\hspace{22mm}\mbox{\guillemotleft{}\rule{1cm}{0.1pt}\guillemotright{}\rule{2cm}{0.1pt}20\rule{0.6cm}{0.1pt}}

\vspace{4mm}
\noindent Diploma project submission date
\hspace{38mm}\mbox{\guillemotleft{}\rule{1cm}{0.1pt}\guillemotright{}\rule{2cm}{0.1pt}20\rule{0.6cm}{0.1pt}}

\vspace{4mm}
\noindent Diploma project initial data

\vspace{2mm}
\noindent\underline{Physical and geological parameters of rocks and media, including density,}\\[1mm]
\noindent\underline{elastic properties, thermal expansion, thermal conductivity, and heat capacity.}

\vspace{4mm}
\vspace{4mm}
\noindent Details of computations and explanations (list of issues due to be addressed)

\vspace{2mm}
\noindent\underline{1. Literature review on geological media, thermoelasticity, and elastic wave}\\[1mm]
\noindent\underline{propagation}\\[1mm]
\noindent\underline{2. Analysis of physical and thermal properties of rocks}\\[1mm]
\noindent\underline{3. Mathematical formulation of the thermoelastic wave propagation problem}\\[1mm]
\noindent\underline{4. Analysis of directional behavior in geological media}\\[1mm]
\noindent\underline{5. Preparation of visual materials and interpretation of obtained results}

\vspace{4mm}
\vspace{4mm}
\noindent CD containing the digital version of diploma paper and attachments

\vspace{2mm}
\noindent\rule{\textwidth}{0.1pt}

\clearpage
\thispagestyle{empty}

\vspace{4mm}
\noindent Consultations on diploma project (with related project chapters named)
\vspace{4mm}

\vspace{2mm}
\noindent\begingroup
\setlength{\tabcolsep}{4pt}
\renewcommand{\arraystretch}{1.15}
\begin{tabular}{|>{\centering\arraybackslash}m{3.1cm}|>{\centering\arraybackslash}m{6.0cm}|>{\centering\arraybackslash}m{3.0cm}|>{\centering\arraybackslash}m{3.0cm}|}
\hline
Consultant & Name & \multicolumn{2}{>{\centering\arraybackslash}m{6.15cm}|}{\rule{0pt}{4mm}Signature, date} \\ \cline{3-4}
& & \begin{tabular}{@{}c@{}}Assignment\\given\end{tabular} & \begin{tabular}{@{}c@{}}Assignment\\received\end{tabular} \\ \hline
\begin{tabular}{@{}c@{}}Compliance\\monitor\end{tabular} & \begin{tabular}{@{}c@{}}M.S., Senior Lecturer,\\Srazh A.S.\end{tabular} & & \\[7mm] \hline
\end{tabular}
\endgroup

\vspace{8mm}
\noindent Date \guillemotleft{}\rule{1cm}{0.1pt}\guillemotright{}\rule{2cm}{0.1pt}20\rule{0.6cm}{0.1pt}

\vspace{4mm}
\noindent Research advisor \hfill \rule{7cm}{0.1pt}

\hfill \small{(signature)}

\vspace{2mm}
\noindent Assignment received by \hfill \rule{7cm}{0.1pt}

\hfill \small{(signature)}

\newpage
\thispagestyle{empty}
\enlargethispage{10mm}

% ---------------------------------------------------------
% DIPLOMA PROJECT WRITING SCHEDULE
% ---------------------------------------------------------
\begin{center}
{Diploma project writing schedule}

\vspace{3mm}
\textbf{Nurshanov D.S., Askarov I.Y., Kassymkhanov T.G., Kaidrakhmetova D.Y.}

\vspace{3mm}
\noindent Title: \textbf{AI-Guided Prediction of Thermoelastic Wave Propagation in Geological Media}
\end{center}

\begin{table}[h!]
\centering
\small
\setlength{\tabcolsep}{4pt}
\renewcommand{\arraystretch}{1.15}

\makebox[\textwidth][c]{%
\begin{tabular}{@{}|>{\centering\arraybackslash}m{0.9cm}|m{11.2cm}|>{\centering\arraybackslash}m{3.4cm}|@{}}
\hline
No. & \centering\arraybackslash Assignment & Period \\ \hline
1 & Creation of the graduation paper writing schedule & September 15 \\ \hline
2 & Collection and analysis of theoretical and geological information & November \\ \hline
3 & Study of thermoelasticity, heat conduction, and elastic wave propagation theory (Introduction, Chapters 1-4) & December \\ \hline
4 & Mathematical formulation of the thermoelastic wave propagation problem & January \\ \hline
5 & Analysis of Existing Models and Approaches (Chapter 5) & January -- February \\ \hline
6 & Development and Implementation of Predictive Models & March -- April \\ \hline
7 & Preparation of visual materials, analysis, and interpretation (Chapters 6, 7) & May \\ \hline
8 & First pre-defense & April 11 \\ \hline
9 & Second pre-defense & May 18 \\ \hline
10 & Submission of the diploma paper for plagiarism check & May 11--15 \\ \hline
11 & Submission of the diploma paper to the compliance monitor & May 16--25 \\ \hline
12 & Submission to the reviewer for approval & May 28 \\ \hline
13 & Diploma project defense & June 1--5 \\ \hline
\end{tabular}%
}
\end{table}
\vspace{1mm}

\noindent\begingroup
\setlength{\tabcolsep}{0pt}
\renewcommand{\arraystretch}{1.00}
\begin{tabular}{@{}p{10.4cm}>{\centering\arraybackslash}p{4.8cm}@{}}
Student: Nurshanov D.S. & \rule{4cm}{0.1pt} \\
& \small{(signature)} \\
Student: Askarov I.Y. & \rule{4cm}{0.1pt} \\
& \small{(signature)} \\
Student: Kassymkhanov T.G. & \rule{4cm}{0.1pt} \\
& \small{(signature)} \\
Student: Kaidrakhmetova D.Y. & \rule{4cm}{0.1pt} \\
& \small{(signature)} \\
Research advisor: Alipova B. N. & \rule{4cm}{0.1pt} \\
& \small{(signature)}
\end{tabular}
\endgroup

\vspace{1mm}
\noindent Date \guillemotleft{}\rule{1cm}{0.1pt}\guillemotright{}\ \rule{1.5cm}{0.1pt}\ 20\rule{0.5cm}{0.1pt}

% % annotations.tex
% This file should be included after the title pages:
% \thispagestyle{empty}

% ---------------------------------------------------------
% TITLE PAGE
% ---------------------------------------------------------
\begin{center}

{MINISTRY OF SCIENCE AND HIGHER EDUCATION OF THE REPUBLIC OF KAZAKHSTAN}\\
INTERNATIONAL INFORMATION TECHNOLOGY UNIVERSITY\\
FACULTY OF COMPUTER TECHNOLOGIES AND CYBERSECURITY

\vspace{35mm}

Nurshanov D.S.\\
Askarov I.Y.\\
Kassymkhanov T.G.\\
Kaidrakhmetova D.Y.

\vspace{15mm}
\textbf{AI-Guided Prediction of Thermoelastic Wave Propagation in Geological Media}

\vspace{10mm}
{DIPLOMA PROJECT}

\vspace{13mm}
\textbf{Educational Program: 6B06112  Data Science}

\vfill

{Almaty 2026}

\end{center}

\newpage
\thispagestyle{empty}

% ---------------------------------------------------------
% APPROVAL PAGE
% ---------------------------------------------------------
\begin{center}
{MINISTRY OF SCIENCE AND HIGHER EDUCATION OF THE REPUBLIC OF KAZAKHSTAN}\\
INTERNATIONAL INFORMATION TECHNOLOGY UNIVERSITY\\
FACULTY OF COMPUTER TECHNOLOGIES AND CYBERSECURITY\\
DEPARTMENT OF MATHEMATICAL AND COMPUTER MODELING
\end{center}

\vspace{4mm}

\begin{flushright}
\begin{minipage}{6.2cm}
Approved\\
Head of Department,\\
PhD, Associate Professor\\
Abdikalikova Z. T.\\[1mm]
\mbox{\guillemotleft{}\rule{0.8cm}{0.1pt}\guillemotright{}\ \rule{2.2cm}{0.1pt}\ 20\rule{0.4cm}{0.1pt}}
\end{minipage}
\end{flushright}

\vspace{8mm}

\begin{center}
\textbf{DIPLOMA PROJECT}

\vspace{4mm}
\textbf{AI-Guided Prediction of Thermoelastic Wave Propagation in Geological Media}

\vspace{5mm}
\textbf{Educational Program: 6B06112 Data Science}
\end{center}

\vspace{8mm}

\noindent
\begingroup
\setlength{\tabcolsep}{0pt}
\renewcommand{\arraystretch}{1.0}
\begin{tabular*}{\textwidth}{@{\extracolsep{\fill}}l c l@{}}
Done by students: & & \\
Nurshanov D.S. & \rule{3.7cm}{0.1pt} & \mbox{\guillemotleft{}\rule{0.8cm}{0.1pt}\guillemotright{}\ \rule{2.2cm}{0.1pt}\ 20\rule{0.4cm}{0.1pt}} \\
& \small{(signature)} & \\[2mm]
Askarov I.Y. & \rule{3.7cm}{0.1pt} & \mbox{\guillemotleft{}\rule{0.8cm}{0.1pt}\guillemotright{}\ \rule{2.2cm}{0.1pt}\ 20\rule{0.4cm}{0.1pt}} \\
& \small{(signature)} & \\[2mm]
Kassymkhanov T.G. & \rule{3.7cm}{0.1pt} & \mbox{\guillemotleft{}\rule{0.8cm}{0.1pt}\guillemotright{}\ \rule{2.2cm}{0.1pt}\ 20\rule{0.4cm}{0.1pt}} \\
& \small{(signature)} & \\[2mm]
Kaidrakhmetova D.Y. & \rule{3.7cm}{0.1pt} & \mbox{\guillemotleft{}\rule{0.8cm}{0.1pt}\guillemotright{}\ \rule{2.2cm}{0.1pt}\ 20\rule{0.4cm}{0.1pt}} \\
& \small{(signature)} & \\[4mm]
\begin{tabular}[t]{@{}l@{}}
Research advisor:\\
Associate Professor,\\
Candidate of Physical and\\
Mathematical Sciences,\\
PhD in Applied Mathematics\\
Alipova B. N.
\end{tabular}
&
\begin{tabular}[t]{@{}c@{}}
\rule{3.7cm}{0.1pt}\\[-1.5mm]
{\small(signature)}
\end{tabular}
&
\mbox{\guillemotleft{}\rule{0.8cm}{0.1pt}\guillemotright{}\ \rule{2.2cm}{0.1pt}\ 20\rule{0.4cm}{0.1pt}}
\end{tabular*}
\endgroup

\newpage
\thispagestyle{empty}
\enlargethispage{10mm}

% ---------------------------------------------------------
% DIPLOMA PROJECT ASSIGNMENT
% ---------------------------------------------------------
\begin{center}
International Information Technology University\\
Faculty of Computer Technologies and Cybersecurity\\
Department of Mathematical and Computer Modeling
\end{center}

\vspace{9mm}

\noindent Educational Program: 6B06112 -- Data Science

\vspace{4mm}
\noindent Diploma Project Assignment

\vspace{4mm}
\noindent Students\\
\textbf{Nurshanov D.S., Askarov I.Y., Kassymkhanov T.G., Kaidrakhmetova D.Y.}

\vspace{4mm}
\noindent Diploma project topic\\
\textbf{AI-Guided Prediction of Thermoelastic Wave Propagation in Geological Media}

\vspace{4mm}
\noindent Approved by IITU order No. \rule{1.4cm}{0.1pt}\ dated
\hspace{22mm}\mbox{\guillemotleft{}\rule{1cm}{0.1pt}\guillemotright{}\rule{2cm}{0.1pt}20\rule{0.6cm}{0.1pt}}

\vspace{4mm}
\noindent Diploma project submission date
\hspace{38mm}\mbox{\guillemotleft{}\rule{1cm}{0.1pt}\guillemotright{}\rule{2cm}{0.1pt}20\rule{0.6cm}{0.1pt}}

\vspace{4mm}
\noindent Diploma project initial data

\vspace{2mm}
\noindent\underline{Physical and geological parameters of rocks and media, including density,}\\[1mm]
\noindent\underline{elastic properties, thermal expansion, thermal conductivity, and heat capacity.}

\vspace{4mm}
\vspace{4mm}
\noindent Details of computations and explanations (list of issues due to be addressed)

\vspace{2mm}
\noindent\underline{1. Literature review on geological media, thermoelasticity, and elastic wave}\\[1mm]
\noindent\underline{propagation}\\[1mm]
\noindent\underline{2. Analysis of physical and thermal properties of rocks}\\[1mm]
\noindent\underline{3. Mathematical formulation of the thermoelastic wave propagation problem}\\[1mm]
\noindent\underline{4. Analysis of directional behavior in geological media}\\[1mm]
\noindent\underline{5. Preparation of visual materials and interpretation of obtained results}

\vspace{4mm}
\vspace{4mm}
\noindent CD containing the digital version of diploma paper and attachments

\vspace{2mm}
\noindent\rule{\textwidth}{0.1pt}

\clearpage
\thispagestyle{empty}

\vspace{4mm}
\noindent Consultations on diploma project (with related project chapters named)
\vspace{4mm}

\vspace{2mm}
\noindent\begingroup
\setlength{\tabcolsep}{4pt}
\renewcommand{\arraystretch}{1.15}
\begin{tabular}{|>{\centering\arraybackslash}m{3.1cm}|>{\centering\arraybackslash}m{6.0cm}|>{\centering\arraybackslash}m{3.0cm}|>{\centering\arraybackslash}m{3.0cm}|}
\hline
Consultant & Name & \multicolumn{2}{>{\centering\arraybackslash}m{6.15cm}|}{\rule{0pt}{4mm}Signature, date} \\ \cline{3-4}
& & \begin{tabular}{@{}c@{}}Assignment\\given\end{tabular} & \begin{tabular}{@{}c@{}}Assignment\\received\end{tabular} \\ \hline
\begin{tabular}{@{}c@{}}Compliance\\monitor\end{tabular} & \begin{tabular}{@{}c@{}}M.S., Senior Lecturer,\\Srazh A.S.\end{tabular} & & \\[7mm] \hline
\end{tabular}
\endgroup

\vspace{8mm}
\noindent Date \guillemotleft{}\rule{1cm}{0.1pt}\guillemotright{}\rule{2cm}{0.1pt}20\rule{0.6cm}{0.1pt}

\vspace{4mm}
\noindent Research advisor \hfill \rule{7cm}{0.1pt}

\hfill \small{(signature)}

\vspace{2mm}
\noindent Assignment received by \hfill \rule{7cm}{0.1pt}

\hfill \small{(signature)}

\newpage
\thispagestyle{empty}
\enlargethispage{10mm}

% ---------------------------------------------------------
% DIPLOMA PROJECT WRITING SCHEDULE
% ---------------------------------------------------------
\begin{center}
{Diploma project writing schedule}

\vspace{3mm}
\textbf{Nurshanov D.S., Askarov I.Y., Kassymkhanov T.G., Kaidrakhmetova D.Y.}

\vspace{3mm}
\noindent Title: \textbf{AI-Guided Prediction of Thermoelastic Wave Propagation in Geological Media}
\end{center}

\begin{table}[h!]
\centering
\small
\setlength{\tabcolsep}{4pt}
\renewcommand{\arraystretch}{1.15}

\makebox[\textwidth][c]{%
\begin{tabular}{@{}|>{\centering\arraybackslash}m{0.9cm}|m{11.2cm}|>{\centering\arraybackslash}m{3.4cm}|@{}}
\hline
No. & \centering\arraybackslash Assignment & Period \\ \hline
1 & Creation of the graduation paper writing schedule & September 15 \\ \hline
2 & Collection and analysis of theoretical and geological information & November \\ \hline
3 & Study of thermoelasticity, heat conduction, and elastic wave propagation theory (Introduction, Chapters 1-4) & December \\ \hline
4 & Mathematical formulation of the thermoelastic wave propagation problem & January \\ \hline
5 & Analysis of Existing Models and Approaches (Chapter 5) & January -- February \\ \hline
6 & Development and Implementation of Predictive Models & March -- April \\ \hline
7 & Preparation of visual materials, analysis, and interpretation (Chapters 6, 7) & May \\ \hline
8 & First pre-defense & April 11 \\ \hline
9 & Second pre-defense & May 18 \\ \hline
10 & Submission of the diploma paper for plagiarism check & May 11--15 \\ \hline
11 & Submission of the diploma paper to the compliance monitor & May 16--25 \\ \hline
12 & Submission to the reviewer for approval & May 28 \\ \hline
13 & Diploma project defense & June 1--5 \\ \hline
\end{tabular}%
}
\end{table}
\vspace{1mm}

\noindent\begingroup
\setlength{\tabcolsep}{0pt}
\renewcommand{\arraystretch}{1.00}
\begin{tabular}{@{}p{10.4cm}>{\centering\arraybackslash}p{4.8cm}@{}}
Student: Nurshanov D.S. & \rule{4cm}{0.1pt} \\
& \small{(signature)} \\
Student: Askarov I.Y. & \rule{4cm}{0.1pt} \\
& \small{(signature)} \\
Student: Kassymkhanov T.G. & \rule{4cm}{0.1pt} \\
& \small{(signature)} \\
Student: Kaidrakhmetova D.Y. & \rule{4cm}{0.1pt} \\
& \small{(signature)} \\
Research advisor: Alipova B. N. & \rule{4cm}{0.1pt} \\
& \small{(signature)}
\end{tabular}
\endgroup

\vspace{1mm}
\noindent Date \guillemotleft{}\rule{1cm}{0.1pt}\guillemotright{}\ \rule{1.5cm}{0.1pt}\ 20\rule{0.5cm}{0.1pt}

% % annotations.tex
% This file should be included after the title pages:
% \thispagestyle{empty}

% ---------------------------------------------------------
% TITLE PAGE
% ---------------------------------------------------------
\begin{center}

{MINISTRY OF SCIENCE AND HIGHER EDUCATION OF THE REPUBLIC OF KAZAKHSTAN}\\
INTERNATIONAL INFORMATION TECHNOLOGY UNIVERSITY\\
FACULTY OF COMPUTER TECHNOLOGIES AND CYBERSECURITY

\vspace{35mm}

Nurshanov D.S.\\
Askarov I.Y.\\
Kassymkhanov T.G.\\
Kaidrakhmetova D.Y.

\vspace{15mm}
\textbf{AI-Guided Prediction of Thermoelastic Wave Propagation in Geological Media}

\vspace{10mm}
{DIPLOMA PROJECT}

\vspace{13mm}
\textbf{Educational Program: 6B06112  Data Science}

\vfill

{Almaty 2026}

\end{center}

\newpage
\thispagestyle{empty}

% ---------------------------------------------------------
% APPROVAL PAGE
% ---------------------------------------------------------
\begin{center}
{MINISTRY OF SCIENCE AND HIGHER EDUCATION OF THE REPUBLIC OF KAZAKHSTAN}\\
INTERNATIONAL INFORMATION TECHNOLOGY UNIVERSITY\\
FACULTY OF COMPUTER TECHNOLOGIES AND CYBERSECURITY\\
DEPARTMENT OF MATHEMATICAL AND COMPUTER MODELING
\end{center}

\vspace{4mm}

\begin{flushright}
\begin{minipage}{6.2cm}
Approved\\
Head of Department,\\
PhD, Associate Professor\\
Abdikalikova Z. T.\\[1mm]
\mbox{\guillemotleft{}\rule{0.8cm}{0.1pt}\guillemotright{}\ \rule{2.2cm}{0.1pt}\ 20\rule{0.4cm}{0.1pt}}
\end{minipage}
\end{flushright}

\vspace{8mm}

\begin{center}
\textbf{DIPLOMA PROJECT}

\vspace{4mm}
\textbf{AI-Guided Prediction of Thermoelastic Wave Propagation in Geological Media}

\vspace{5mm}
\textbf{Educational Program: 6B06112 Data Science}
\end{center}

\vspace{8mm}

\noindent
\begingroup
\setlength{\tabcolsep}{0pt}
\renewcommand{\arraystretch}{1.0}
\begin{tabular*}{\textwidth}{@{\extracolsep{\fill}}l c l@{}}
Done by students: & & \\
Nurshanov D.S. & \rule{3.7cm}{0.1pt} & \mbox{\guillemotleft{}\rule{0.8cm}{0.1pt}\guillemotright{}\ \rule{2.2cm}{0.1pt}\ 20\rule{0.4cm}{0.1pt}} \\
& \small{(signature)} & \\[2mm]
Askarov I.Y. & \rule{3.7cm}{0.1pt} & \mbox{\guillemotleft{}\rule{0.8cm}{0.1pt}\guillemotright{}\ \rule{2.2cm}{0.1pt}\ 20\rule{0.4cm}{0.1pt}} \\
& \small{(signature)} & \\[2mm]
Kassymkhanov T.G. & \rule{3.7cm}{0.1pt} & \mbox{\guillemotleft{}\rule{0.8cm}{0.1pt}\guillemotright{}\ \rule{2.2cm}{0.1pt}\ 20\rule{0.4cm}{0.1pt}} \\
& \small{(signature)} & \\[2mm]
Kaidrakhmetova D.Y. & \rule{3.7cm}{0.1pt} & \mbox{\guillemotleft{}\rule{0.8cm}{0.1pt}\guillemotright{}\ \rule{2.2cm}{0.1pt}\ 20\rule{0.4cm}{0.1pt}} \\
& \small{(signature)} & \\[4mm]
\begin{tabular}[t]{@{}l@{}}
Research advisor:\\
Associate Professor,\\
Candidate of Physical and\\
Mathematical Sciences,\\
PhD in Applied Mathematics\\
Alipova B. N.
\end{tabular}
&
\begin{tabular}[t]{@{}c@{}}
\rule{3.7cm}{0.1pt}\\[-1.5mm]
{\small(signature)}
\end{tabular}
&
\mbox{\guillemotleft{}\rule{0.8cm}{0.1pt}\guillemotright{}\ \rule{2.2cm}{0.1pt}\ 20\rule{0.4cm}{0.1pt}}
\end{tabular*}
\endgroup

\newpage
\thispagestyle{empty}
\enlargethispage{10mm}

% ---------------------------------------------------------
% DIPLOMA PROJECT ASSIGNMENT
% ---------------------------------------------------------
\begin{center}
International Information Technology University\\
Faculty of Computer Technologies and Cybersecurity\\
Department of Mathematical and Computer Modeling
\end{center}

\vspace{9mm}

\noindent Educational Program: 6B06112 -- Data Science

\vspace{4mm}
\noindent Diploma Project Assignment

\vspace{4mm}
\noindent Students\\
\textbf{Nurshanov D.S., Askarov I.Y., Kassymkhanov T.G., Kaidrakhmetova D.Y.}

\vspace{4mm}
\noindent Diploma project topic\\
\textbf{AI-Guided Prediction of Thermoelastic Wave Propagation in Geological Media}

\vspace{4mm}
\noindent Approved by IITU order No. \rule{1.4cm}{0.1pt}\ dated
\hspace{22mm}\mbox{\guillemotleft{}\rule{1cm}{0.1pt}\guillemotright{}\rule{2cm}{0.1pt}20\rule{0.6cm}{0.1pt}}

\vspace{4mm}
\noindent Diploma project submission date
\hspace{38mm}\mbox{\guillemotleft{}\rule{1cm}{0.1pt}\guillemotright{}\rule{2cm}{0.1pt}20\rule{0.6cm}{0.1pt}}

\vspace{4mm}
\noindent Diploma project initial data

\vspace{2mm}
\noindent\underline{Physical and geological parameters of rocks and media, including density,}\\[1mm]
\noindent\underline{elastic properties, thermal expansion, thermal conductivity, and heat capacity.}

\vspace{4mm}
\vspace{4mm}
\noindent Details of computations and explanations (list of issues due to be addressed)

\vspace{2mm}
\noindent\underline{1. Literature review on geological media, thermoelasticity, and elastic wave}\\[1mm]
\noindent\underline{propagation}\\[1mm]
\noindent\underline{2. Analysis of physical and thermal properties of rocks}\\[1mm]
\noindent\underline{3. Mathematical formulation of the thermoelastic wave propagation problem}\\[1mm]
\noindent\underline{4. Analysis of directional behavior in geological media}\\[1mm]
\noindent\underline{5. Preparation of visual materials and interpretation of obtained results}

\vspace{4mm}
\vspace{4mm}
\noindent CD containing the digital version of diploma paper and attachments

\vspace{2mm}
\noindent\rule{\textwidth}{0.1pt}

\clearpage
\thispagestyle{empty}

\vspace{4mm}
\noindent Consultations on diploma project (with related project chapters named)
\vspace{4mm}

\vspace{2mm}
\noindent\begingroup
\setlength{\tabcolsep}{4pt}
\renewcommand{\arraystretch}{1.15}
\begin{tabular}{|>{\centering\arraybackslash}m{3.1cm}|>{\centering\arraybackslash}m{6.0cm}|>{\centering\arraybackslash}m{3.0cm}|>{\centering\arraybackslash}m{3.0cm}|}
\hline
Consultant & Name & \multicolumn{2}{>{\centering\arraybackslash}m{6.15cm}|}{\rule{0pt}{4mm}Signature, date} \\ \cline{3-4}
& & \begin{tabular}{@{}c@{}}Assignment\\given\end{tabular} & \begin{tabular}{@{}c@{}}Assignment\\received\end{tabular} \\ \hline
\begin{tabular}{@{}c@{}}Compliance\\monitor\end{tabular} & \begin{tabular}{@{}c@{}}M.S., Senior Lecturer,\\Srazh A.S.\end{tabular} & & \\[7mm] \hline
\end{tabular}
\endgroup

\vspace{8mm}
\noindent Date \guillemotleft{}\rule{1cm}{0.1pt}\guillemotright{}\rule{2cm}{0.1pt}20\rule{0.6cm}{0.1pt}

\vspace{4mm}
\noindent Research advisor \hfill \rule{7cm}{0.1pt}

\hfill \small{(signature)}

\vspace{2mm}
\noindent Assignment received by \hfill \rule{7cm}{0.1pt}

\hfill \small{(signature)}

\newpage
\thispagestyle{empty}
\enlargethispage{10mm}

% ---------------------------------------------------------
% DIPLOMA PROJECT WRITING SCHEDULE
% ---------------------------------------------------------
\begin{center}
{Diploma project writing schedule}

\vspace{3mm}
\textbf{Nurshanov D.S., Askarov I.Y., Kassymkhanov T.G., Kaidrakhmetova D.Y.}

\vspace{3mm}
\noindent Title: \textbf{AI-Guided Prediction of Thermoelastic Wave Propagation in Geological Media}
\end{center}

\begin{table}[h!]
\centering
\small
\setlength{\tabcolsep}{4pt}
\renewcommand{\arraystretch}{1.15}

\makebox[\textwidth][c]{%
\begin{tabular}{@{}|>{\centering\arraybackslash}m{0.9cm}|m{11.2cm}|>{\centering\arraybackslash}m{3.4cm}|@{}}
\hline
No. & \centering\arraybackslash Assignment & Period \\ \hline
1 & Creation of the graduation paper writing schedule & September 15 \\ \hline
2 & Collection and analysis of theoretical and geological information & November \\ \hline
3 & Study of thermoelasticity, heat conduction, and elastic wave propagation theory (Introduction, Chapters 1-4) & December \\ \hline
4 & Mathematical formulation of the thermoelastic wave propagation problem & January \\ \hline
5 & Analysis of Existing Models and Approaches (Chapter 5) & January -- February \\ \hline
6 & Development and Implementation of Predictive Models & March -- April \\ \hline
7 & Preparation of visual materials, analysis, and interpretation (Chapters 6, 7) & May \\ \hline
8 & First pre-defense & April 11 \\ \hline
9 & Second pre-defense & May 18 \\ \hline
10 & Submission of the diploma paper for plagiarism check & May 11--15 \\ \hline
11 & Submission of the diploma paper to the compliance monitor & May 16--25 \\ \hline
12 & Submission to the reviewer for approval & May 28 \\ \hline
13 & Diploma project defense & June 1--5 \\ \hline
\end{tabular}%
}
\end{table}
\vspace{1mm}

\noindent\begingroup
\setlength{\tabcolsep}{0pt}
\renewcommand{\arraystretch}{1.00}
\begin{tabular}{@{}p{10.4cm}>{\centering\arraybackslash}p{4.8cm}@{}}
Student: Nurshanov D.S. & \rule{4cm}{0.1pt} \\
& \small{(signature)} \\
Student: Askarov I.Y. & \rule{4cm}{0.1pt} \\
& \small{(signature)} \\
Student: Kassymkhanov T.G. & \rule{4cm}{0.1pt} \\
& \small{(signature)} \\
Student: Kaidrakhmetova D.Y. & \rule{4cm}{0.1pt} \\
& \small{(signature)} \\
Research advisor: Alipova B. N. & \rule{4cm}{0.1pt} \\
& \small{(signature)}
\end{tabular}
\endgroup

\vspace{1mm}
\noindent Date \guillemotleft{}\rule{1cm}{0.1pt}\guillemotright{}\ \rule{1.5cm}{0.1pt}\ 20\rule{0.5cm}{0.1pt}

% % annotations.tex
% This file should be included after the title pages:
% \input{tiles.tex}
% \input{annotations.tex}

\clearpage
\thispagestyle{empty}

\begin{center}
\textbf{\Large АҢДАТПА}
\end{center}

Дипломдық жобаның мақсаты геологиялық ортада термосерпімді толқындардың таралу бағытын болжауға арналған теориялық және есептеу негізін әзірлеу болып табылады. Жұмыста тау жыныстары табиғи қатты денелер ретінде қарастырылып, олардың физикалық мінез-құлқы минералдық құрамға, тығыздыққа, кеуектілікке, серпімділік қасиеттеріне, жылу өткізгіштікке, жылу сыйымдылығына, термиялық ұлғаю коэффициентіне, температуралық жағдайға және кернеулік күйге тәуелді екені сипатталады.

Жобаның теориялық бөлігі термосерпімділік, жылу өткізгіштік, термиялық кернеу, серпімді деформация және геологиялық ортадағы толқындардың таралуы арасындағы байланысты түсіндіруге бағытталған. Ерекше назар қабаттылық, жарықшақтар, кеуектер, анизотропия және жыныстың ішкі құрылымы толқындардың бағытқа тәуелді таралуына қалай әсер ететініне аударылады. Температураның өзгеруі еркін ұлғаю немесе сығылу ғана емес, геологиялық шектеулер жағдайында қосымша термиялық кернеулер тудыруы мүмкін екені көрсетіледі.

Практикалық бөлік геологиялық материалдардың физикалық параметрлеріне сүйене отырып, термосерпімді толқындардың таралу бағытын салыстырмалы түрде талдауға арналған прототиптік тәсілді қарастырады. Бұл жұмыс нақты далалық геофизикалық модельдеуді толық алмастыруды мақсат етпейді; оның негізгі міндеті — геологиялық ортадағы жылулық және механикалық процестердің байланысын теориялық тұрғыдан негіздеп, бағыттық толқындық мінез-құлықты түсіндіруге арналған құрылымдалған негіз ұсыну.

Дипломдық жобада \rule{1cm}{0.1pt} бет, \rule{1cm}{0.1pt} кесте, \rule{1cm}{0.1pt} сурет және \rule{1cm}{0.1pt} пайдаланылған әдебиеттер бар.

\vspace{5mm}
\noindent\textbf{Кілт сөздер тізімі:} ТЕРМОСЕРПІМДІЛІК, ГЕОЛОГИЯЛЫҚ ОРТА, СЕРПІМДІ ТОЛҚЫНДАР, ЖЫЛУ ӨТКІЗГІШТІК, ТЕРМИЯЛЫҚ КЕРНЕУ, ТАУ ЖЫНЫСТАРЫ, БАҒЫТТЫҚ ТАРАЛУ, АНИЗОТРОПИЯ, ГЕОФИЗИКА, ЖАСАНДЫ ИНТЕЛЛЕКТ.

\clearpage
\thispagestyle{empty}

\begin{center}
\textbf{\Large АННОТАЦИЯ}
\end{center}

Целью дипломного проекта является разработка и теоретическое обоснование подхода к направленному прогнозированию распространения термоупругих волн в геологических средах. В работе геологические среды рассматриваются как природные твердые материалы, физическое поведение которых зависит от минерального состава, плотности, пористости, упругих свойств, теплопроводности, теплоемкости, коэффициента теплового расширения, температурного режима и напряженного состояния.

Теоретическая часть проекта направлена на объяснение связи между термоупругостью, теплопроводностью, тепловыми напряжениями, упругой деформацией и распространением волн в горных породах. Особое внимание уделяется тому, как слоистость, трещиноватость, поровая структура, анизотропия и внутренняя структура пород влияют на направленное поведение волн. Показано, что изменение температуры может приводить не только к свободному расширению или сжатию, но и к возникновению дополнительных тепловых напряжений при наличии геологических ограничений.

Практическая часть работы рассматривает прототипный подход к сравнительному анализу направления распространения термоупругих волн на основе физических параметров геологических материалов. Работа не претендует на полную замену реального полевого геофизического моделирования; ее основная задача состоит в теоретическом обосновании связи тепловых и механических процессов в геологической среде и в формировании структурированной основы для интерпретации направленного волнового поведения.

Дипломный проект содержит \rule{1cm}{0.1pt} страниц, \rule{1cm}{0.1pt} таблиц, \rule{1cm}{0.1pt} иллюстраций и \rule{1cm}{0.1pt} использованных источников.

\vspace{5mm}
\noindent\textbf{Перечень ключевых слов:} ТЕРМОУПРУГОСТЬ, ГЕОЛОГИЧЕСКАЯ СРЕДА, УПРУГИЕ ВОЛНЫ, ТЕПЛОПРОВОДНОСТЬ, ТЕПЛОВОЕ НАПРЯЖЕНИЕ, ГОРНЫЕ ПОРОДЫ, НАПРАВЛЕННОЕ РАСПРОСТРАНЕНИЕ, АНИЗОТРОПИЯ, ГЕОФИЗИКА, ИСКУССТВЕННЫЙ ИНТЕЛЛЕКТ.

\clearpage
\thispagestyle{empty}

\begin{center}
\textbf{\Large ABSTRACT}
\end{center}

The aim of the diploma project is to develop and theoretically justify an approach to directional prediction of thermoelastic wave propagation in geological media. Geological media are considered as natural solid materials whose physical behavior depends on mineral composition, density, porosity, elastic properties, thermal conductivity, heat capacity, coefficient of thermal expansion, temperature conditions, and stress state.

The theoretical part of the project explains the relationship between thermoelasticity, heat conduction, thermal stress, elastic deformation, and wave propagation in rocks. Particular attention is given to the influence of layering, fractures, pore structure, anisotropy, and internal rock fabric on directional wave behavior. The work emphasizes that temperature changes may cause not only free expansion or contraction, but also additional thermal stresses when geological constraints are present.

The practical part considers a prototype-oriented approach for comparative analysis of the direction of thermoelastic wave propagation based on physical parameters of geological materials. The work does not claim to replace full-scale field geophysical modelling. Its main purpose is to provide a theoretical explanation of the interaction between thermal and mechanical processes in geological media and to form a structured basis for interpreting directional wave behavior.

The graduation project contains \rule{1cm}{0.1pt} pages, \rule{1cm}{0.1pt} tables, \rule{1cm}{0.1pt} illustrations, and \rule{1cm}{0.1pt} references.

\vspace{5mm}
\noindent\textbf{Keywords:} THERMOELASTICITY, GEOLOGICAL MEDIA, ELASTIC WAVES, HEAT CONDUCTION, THERMAL STRESS, ROCKS, DIRECTIONAL PROPAGATION, ANISOTROPY, GEOPHYSICS, ARTIFICIAL INTELLIGENCE.


\clearpage
\thispagestyle{empty}

\begin{center}
\textbf{\Large АҢДАТПА}
\end{center}

Дипломдық жобаның мақсаты геологиялық ортада термосерпімді толқындардың таралу бағытын болжауға арналған теориялық және есептеу негізін әзірлеу болып табылады. Жұмыста тау жыныстары табиғи қатты денелер ретінде қарастырылып, олардың физикалық мінез-құлқы минералдық құрамға, тығыздыққа, кеуектілікке, серпімділік қасиеттеріне, жылу өткізгіштікке, жылу сыйымдылығына, термиялық ұлғаю коэффициентіне, температуралық жағдайға және кернеулік күйге тәуелді екені сипатталады.

Жобаның теориялық бөлігі термосерпімділік, жылу өткізгіштік, термиялық кернеу, серпімді деформация және геологиялық ортадағы толқындардың таралуы арасындағы байланысты түсіндіруге бағытталған. Ерекше назар қабаттылық, жарықшақтар, кеуектер, анизотропия және жыныстың ішкі құрылымы толқындардың бағытқа тәуелді таралуына қалай әсер ететініне аударылады. Температураның өзгеруі еркін ұлғаю немесе сығылу ғана емес, геологиялық шектеулер жағдайында қосымша термиялық кернеулер тудыруы мүмкін екені көрсетіледі.

Практикалық бөлік геологиялық материалдардың физикалық параметрлеріне сүйене отырып, термосерпімді толқындардың таралу бағытын салыстырмалы түрде талдауға арналған прототиптік тәсілді қарастырады. Бұл жұмыс нақты далалық геофизикалық модельдеуді толық алмастыруды мақсат етпейді; оның негізгі міндеті — геологиялық ортадағы жылулық және механикалық процестердің байланысын теориялық тұрғыдан негіздеп, бағыттық толқындық мінез-құлықты түсіндіруге арналған құрылымдалған негіз ұсыну.

Дипломдық жобада \rule{1cm}{0.1pt} бет, \rule{1cm}{0.1pt} кесте, \rule{1cm}{0.1pt} сурет және \rule{1cm}{0.1pt} пайдаланылған әдебиеттер бар.

\vspace{5mm}
\noindent\textbf{Кілт сөздер тізімі:} ТЕРМОСЕРПІМДІЛІК, ГЕОЛОГИЯЛЫҚ ОРТА, СЕРПІМДІ ТОЛҚЫНДАР, ЖЫЛУ ӨТКІЗГІШТІК, ТЕРМИЯЛЫҚ КЕРНЕУ, ТАУ ЖЫНЫСТАРЫ, БАҒЫТТЫҚ ТАРАЛУ, АНИЗОТРОПИЯ, ГЕОФИЗИКА, ЖАСАНДЫ ИНТЕЛЛЕКТ.

\clearpage
\thispagestyle{empty}

\begin{center}
\textbf{\Large АННОТАЦИЯ}
\end{center}

Целью дипломного проекта является разработка и теоретическое обоснование подхода к направленному прогнозированию распространения термоупругих волн в геологических средах. В работе геологические среды рассматриваются как природные твердые материалы, физическое поведение которых зависит от минерального состава, плотности, пористости, упругих свойств, теплопроводности, теплоемкости, коэффициента теплового расширения, температурного режима и напряженного состояния.

Теоретическая часть проекта направлена на объяснение связи между термоупругостью, теплопроводностью, тепловыми напряжениями, упругой деформацией и распространением волн в горных породах. Особое внимание уделяется тому, как слоистость, трещиноватость, поровая структура, анизотропия и внутренняя структура пород влияют на направленное поведение волн. Показано, что изменение температуры может приводить не только к свободному расширению или сжатию, но и к возникновению дополнительных тепловых напряжений при наличии геологических ограничений.

Практическая часть работы рассматривает прототипный подход к сравнительному анализу направления распространения термоупругих волн на основе физических параметров геологических материалов. Работа не претендует на полную замену реального полевого геофизического моделирования; ее основная задача состоит в теоретическом обосновании связи тепловых и механических процессов в геологической среде и в формировании структурированной основы для интерпретации направленного волнового поведения.

Дипломный проект содержит \rule{1cm}{0.1pt} страниц, \rule{1cm}{0.1pt} таблиц, \rule{1cm}{0.1pt} иллюстраций и \rule{1cm}{0.1pt} использованных источников.

\vspace{5mm}
\noindent\textbf{Перечень ключевых слов:} ТЕРМОУПРУГОСТЬ, ГЕОЛОГИЧЕСКАЯ СРЕДА, УПРУГИЕ ВОЛНЫ, ТЕПЛОПРОВОДНОСТЬ, ТЕПЛОВОЕ НАПРЯЖЕНИЕ, ГОРНЫЕ ПОРОДЫ, НАПРАВЛЕННОЕ РАСПРОСТРАНЕНИЕ, АНИЗОТРОПИЯ, ГЕОФИЗИКА, ИСКУССТВЕННЫЙ ИНТЕЛЛЕКТ.

\clearpage
\thispagestyle{empty}

\begin{center}
\textbf{\Large ABSTRACT}
\end{center}

The aim of the diploma project is to develop and theoretically justify an approach to directional prediction of thermoelastic wave propagation in geological media. Geological media are considered as natural solid materials whose physical behavior depends on mineral composition, density, porosity, elastic properties, thermal conductivity, heat capacity, coefficient of thermal expansion, temperature conditions, and stress state.

The theoretical part of the project explains the relationship between thermoelasticity, heat conduction, thermal stress, elastic deformation, and wave propagation in rocks. Particular attention is given to the influence of layering, fractures, pore structure, anisotropy, and internal rock fabric on directional wave behavior. The work emphasizes that temperature changes may cause not only free expansion or contraction, but also additional thermal stresses when geological constraints are present.

The practical part considers a prototype-oriented approach for comparative analysis of the direction of thermoelastic wave propagation based on physical parameters of geological materials. The work does not claim to replace full-scale field geophysical modelling. Its main purpose is to provide a theoretical explanation of the interaction between thermal and mechanical processes in geological media and to form a structured basis for interpreting directional wave behavior.

The graduation project contains \rule{1cm}{0.1pt} pages, \rule{1cm}{0.1pt} tables, \rule{1cm}{0.1pt} illustrations, and \rule{1cm}{0.1pt} references.

\vspace{5mm}
\noindent\textbf{Keywords:} THERMOELASTICITY, GEOLOGICAL MEDIA, ELASTIC WAVES, HEAT CONDUCTION, THERMAL STRESS, ROCKS, DIRECTIONAL PROPAGATION, ANISOTROPY, GEOPHYSICS, ARTIFICIAL INTELLIGENCE.


\clearpage
\thispagestyle{empty}

\begin{center}
\textbf{\Large АҢДАТПА}
\end{center}

Дипломдық жобаның мақсаты геологиялық ортада термосерпімді толқындардың таралу бағытын болжауға арналған теориялық және есептеу негізін әзірлеу болып табылады. Жұмыста тау жыныстары табиғи қатты денелер ретінде қарастырылып, олардың физикалық мінез-құлқы минералдық құрамға, тығыздыққа, кеуектілікке, серпімділік қасиеттеріне, жылу өткізгіштікке, жылу сыйымдылығына, термиялық ұлғаю коэффициентіне, температуралық жағдайға және кернеулік күйге тәуелді екені сипатталады.

Жобаның теориялық бөлігі термосерпімділік, жылу өткізгіштік, термиялық кернеу, серпімді деформация және геологиялық ортадағы толқындардың таралуы арасындағы байланысты түсіндіруге бағытталған. Ерекше назар қабаттылық, жарықшақтар, кеуектер, анизотропия және жыныстың ішкі құрылымы толқындардың бағытқа тәуелді таралуына қалай әсер ететініне аударылады. Температураның өзгеруі еркін ұлғаю немесе сығылу ғана емес, геологиялық шектеулер жағдайында қосымша термиялық кернеулер тудыруы мүмкін екені көрсетіледі.

Практикалық бөлік геологиялық материалдардың физикалық параметрлеріне сүйене отырып, термосерпімді толқындардың таралу бағытын салыстырмалы түрде талдауға арналған прототиптік тәсілді қарастырады. Бұл жұмыс нақты далалық геофизикалық модельдеуді толық алмастыруды мақсат етпейді; оның негізгі міндеті — геологиялық ортадағы жылулық және механикалық процестердің байланысын теориялық тұрғыдан негіздеп, бағыттық толқындық мінез-құлықты түсіндіруге арналған құрылымдалған негіз ұсыну.

Дипломдық жобада \rule{1cm}{0.1pt} бет, \rule{1cm}{0.1pt} кесте, \rule{1cm}{0.1pt} сурет және \rule{1cm}{0.1pt} пайдаланылған әдебиеттер бар.

\vspace{5mm}
\noindent\textbf{Кілт сөздер тізімі:} ТЕРМОСЕРПІМДІЛІК, ГЕОЛОГИЯЛЫҚ ОРТА, СЕРПІМДІ ТОЛҚЫНДАР, ЖЫЛУ ӨТКІЗГІШТІК, ТЕРМИЯЛЫҚ КЕРНЕУ, ТАУ ЖЫНЫСТАРЫ, БАҒЫТТЫҚ ТАРАЛУ, АНИЗОТРОПИЯ, ГЕОФИЗИКА, ЖАСАНДЫ ИНТЕЛЛЕКТ.

\clearpage
\thispagestyle{empty}

\begin{center}
\textbf{\Large АННОТАЦИЯ}
\end{center}

Целью дипломного проекта является разработка и теоретическое обоснование подхода к направленному прогнозированию распространения термоупругих волн в геологических средах. В работе геологические среды рассматриваются как природные твердые материалы, физическое поведение которых зависит от минерального состава, плотности, пористости, упругих свойств, теплопроводности, теплоемкости, коэффициента теплового расширения, температурного режима и напряженного состояния.

Теоретическая часть проекта направлена на объяснение связи между термоупругостью, теплопроводностью, тепловыми напряжениями, упругой деформацией и распространением волн в горных породах. Особое внимание уделяется тому, как слоистость, трещиноватость, поровая структура, анизотропия и внутренняя структура пород влияют на направленное поведение волн. Показано, что изменение температуры может приводить не только к свободному расширению или сжатию, но и к возникновению дополнительных тепловых напряжений при наличии геологических ограничений.

Практическая часть работы рассматривает прототипный подход к сравнительному анализу направления распространения термоупругих волн на основе физических параметров геологических материалов. Работа не претендует на полную замену реального полевого геофизического моделирования; ее основная задача состоит в теоретическом обосновании связи тепловых и механических процессов в геологической среде и в формировании структурированной основы для интерпретации направленного волнового поведения.

Дипломный проект содержит \rule{1cm}{0.1pt} страниц, \rule{1cm}{0.1pt} таблиц, \rule{1cm}{0.1pt} иллюстраций и \rule{1cm}{0.1pt} использованных источников.

\vspace{5mm}
\noindent\textbf{Перечень ключевых слов:} ТЕРМОУПРУГОСТЬ, ГЕОЛОГИЧЕСКАЯ СРЕДА, УПРУГИЕ ВОЛНЫ, ТЕПЛОПРОВОДНОСТЬ, ТЕПЛОВОЕ НАПРЯЖЕНИЕ, ГОРНЫЕ ПОРОДЫ, НАПРАВЛЕННОЕ РАСПРОСТРАНЕНИЕ, АНИЗОТРОПИЯ, ГЕОФИЗИКА, ИСКУССТВЕННЫЙ ИНТЕЛЛЕКТ.

\clearpage
\thispagestyle{empty}

\begin{center}
\textbf{\Large ABSTRACT}
\end{center}

The aim of the diploma project is to develop and theoretically justify an approach to directional prediction of thermoelastic wave propagation in geological media. Geological media are considered as natural solid materials whose physical behavior depends on mineral composition, density, porosity, elastic properties, thermal conductivity, heat capacity, coefficient of thermal expansion, temperature conditions, and stress state.

The theoretical part of the project explains the relationship between thermoelasticity, heat conduction, thermal stress, elastic deformation, and wave propagation in rocks. Particular attention is given to the influence of layering, fractures, pore structure, anisotropy, and internal rock fabric on directional wave behavior. The work emphasizes that temperature changes may cause not only free expansion or contraction, but also additional thermal stresses when geological constraints are present.

The practical part considers a prototype-oriented approach for comparative analysis of the direction of thermoelastic wave propagation based on physical parameters of geological materials. The work does not claim to replace full-scale field geophysical modelling. Its main purpose is to provide a theoretical explanation of the interaction between thermal and mechanical processes in geological media and to form a structured basis for interpreting directional wave behavior.

The graduation project contains \rule{1cm}{0.1pt} pages, \rule{1cm}{0.1pt} tables, \rule{1cm}{0.1pt} illustrations, and \rule{1cm}{0.1pt} references.

\vspace{5mm}
\noindent\textbf{Keywords:} THERMOELASTICITY, GEOLOGICAL MEDIA, ELASTIC WAVES, HEAT CONDUCTION, THERMAL STRESS, ROCKS, DIRECTIONAL PROPAGATION, ANISOTROPY, GEOPHYSICS, ARTIFICIAL INTELLIGENCE.


\clearpage
\thispagestyle{empty}

\begin{center}
\textbf{\Large АҢДАТПА}
\end{center}

Дипломдық жобаның мақсаты геологиялық ортада термосерпімді толқындардың таралу бағытын болжауға арналған теориялық және есептеу негізін әзірлеу болып табылады. Жұмыста тау жыныстары табиғи қатты денелер ретінде қарастырылып, олардың физикалық мінез-құлқы минералдық құрамға, тығыздыққа, кеуектілікке, серпімділік қасиеттеріне, жылу өткізгіштікке, жылу сыйымдылығына, термиялық ұлғаю коэффициентіне, температуралық жағдайға және кернеулік күйге тәуелді екені сипатталады.

Жобаның теориялық бөлігі термосерпімділік, жылу өткізгіштік, термиялық кернеу, серпімді деформация және геологиялық ортадағы толқындардың таралуы арасындағы байланысты түсіндіруге бағытталған. Ерекше назар қабаттылық, жарықшақтар, кеуектер, анизотропия және жыныстың ішкі құрылымы толқындардың бағытқа тәуелді таралуына қалай әсер ететініне аударылады. Температураның өзгеруі еркін ұлғаю немесе сығылу ғана емес, геологиялық шектеулер жағдайында қосымша термиялық кернеулер тудыруы мүмкін екені көрсетіледі.

Практикалық бөлік геологиялық материалдардың физикалық параметрлеріне сүйене отырып, термосерпімді толқындардың таралу бағытын салыстырмалы түрде талдауға арналған прототиптік тәсілді қарастырады. Бұл жұмыс нақты далалық геофизикалық модельдеуді толық алмастыруды мақсат етпейді; оның негізгі міндеті — геологиялық ортадағы жылулық және механикалық процестердің байланысын теориялық тұрғыдан негіздеп, бағыттық толқындық мінез-құлықты түсіндіруге арналған құрылымдалған негіз ұсыну.

Дипломдық жобада \rule{1cm}{0.1pt} бет, \rule{1cm}{0.1pt} кесте, \rule{1cm}{0.1pt} сурет және \rule{1cm}{0.1pt} пайдаланылған әдебиеттер бар.

\vspace{5mm}
\noindent\textbf{Кілт сөздер тізімі:} ТЕРМОСЕРПІМДІЛІК, ГЕОЛОГИЯЛЫҚ ОРТА, СЕРПІМДІ ТОЛҚЫНДАР, ЖЫЛУ ӨТКІЗГІШТІК, ТЕРМИЯЛЫҚ КЕРНЕУ, ТАУ ЖЫНЫСТАРЫ, БАҒЫТТЫҚ ТАРАЛУ, АНИЗОТРОПИЯ, ГЕОФИЗИКА, ЖАСАНДЫ ИНТЕЛЛЕКТ.

\clearpage
\thispagestyle{empty}

\begin{center}
\textbf{\Large АННОТАЦИЯ}
\end{center}

Целью дипломного проекта является разработка и теоретическое обоснование подхода к направленному прогнозированию распространения термоупругих волн в геологических средах. В работе геологические среды рассматриваются как природные твердые материалы, физическое поведение которых зависит от минерального состава, плотности, пористости, упругих свойств, теплопроводности, теплоемкости, коэффициента теплового расширения, температурного режима и напряженного состояния.

Теоретическая часть проекта направлена на объяснение связи между термоупругостью, теплопроводностью, тепловыми напряжениями, упругой деформацией и распространением волн в горных породах. Особое внимание уделяется тому, как слоистость, трещиноватость, поровая структура, анизотропия и внутренняя структура пород влияют на направленное поведение волн. Показано, что изменение температуры может приводить не только к свободному расширению или сжатию, но и к возникновению дополнительных тепловых напряжений при наличии геологических ограничений.

Практическая часть работы рассматривает прототипный подход к сравнительному анализу направления распространения термоупругих волн на основе физических параметров геологических материалов. Работа не претендует на полную замену реального полевого геофизического моделирования; ее основная задача состоит в теоретическом обосновании связи тепловых и механических процессов в геологической среде и в формировании структурированной основы для интерпретации направленного волнового поведения.

Дипломный проект содержит \rule{1cm}{0.1pt} страниц, \rule{1cm}{0.1pt} таблиц, \rule{1cm}{0.1pt} иллюстраций и \rule{1cm}{0.1pt} использованных источников.

\vspace{5mm}
\noindent\textbf{Перечень ключевых слов:} ТЕРМОУПРУГОСТЬ, ГЕОЛОГИЧЕСКАЯ СРЕДА, УПРУГИЕ ВОЛНЫ, ТЕПЛОПРОВОДНОСТЬ, ТЕПЛОВОЕ НАПРЯЖЕНИЕ, ГОРНЫЕ ПОРОДЫ, НАПРАВЛЕННОЕ РАСПРОСТРАНЕНИЕ, АНИЗОТРОПИЯ, ГЕОФИЗИКА, ИСКУССТВЕННЫЙ ИНТЕЛЛЕКТ.

\clearpage
\thispagestyle{empty}

\begin{center}
\textbf{\Large ABSTRACT}
\end{center}

The aim of the diploma project is to develop and theoretically justify an approach to directional prediction of thermoelastic wave propagation in geological media. Geological media are considered as natural solid materials whose physical behavior depends on mineral composition, density, porosity, elastic properties, thermal conductivity, heat capacity, coefficient of thermal expansion, temperature conditions, and stress state.

The theoretical part of the project explains the relationship between thermoelasticity, heat conduction, thermal stress, elastic deformation, and wave propagation in rocks. Particular attention is given to the influence of layering, fractures, pore structure, anisotropy, and internal rock fabric on directional wave behavior. The work emphasizes that temperature changes may cause not only free expansion or contraction, but also additional thermal stresses when geological constraints are present.

The practical part considers a prototype-oriented approach for comparative analysis of the direction of thermoelastic wave propagation based on physical parameters of geological materials. The work does not claim to replace full-scale field geophysical modelling. Its main purpose is to provide a theoretical explanation of the interaction between thermal and mechanical processes in geological media and to form a structured basis for interpreting directional wave behavior.

The graduation project contains \rule{1cm}{0.1pt} pages, \rule{1cm}{0.1pt} tables, \rule{1cm}{0.1pt} illustrations, and \rule{1cm}{0.1pt} references.

\vspace{5mm}
\noindent\textbf{Keywords:} THERMOELASTICITY, GEOLOGICAL MEDIA, ELASTIC WAVES, HEAT CONDUCTION, THERMAL STRESS, ROCKS, DIRECTIONAL PROPAGATION, ANISOTROPY, GEOPHYSICS, ARTIFICIAL INTELLIGENCE.
